\documentclass[aspectratio=169]{beamer}
\usetheme{Madrid}
\usecolortheme{default}
\usepackage{amsmath}
\usepackage{graphicx}
\usepackage{tikz-cd}
\usepackage{multicol}

\setlength{\parindent}{0pt}
\setlength{\parskip}{1\baselineskip}

\title[Lecture 06]{An Introduction to E-Commerce\\\small with a Focus on Machine Learning, Deep Learning, and Artificial Intelligence}
\subtitle{Lecture 06: ERP/CRM/SCM Integration in Commerce}
\author{Dr. Haitham A. El-Ghareeb}
\institute{Faculty of Computers and Information Sciences \\ Mansoura University \\ Egypt}
\date{\today}

\begin{document}

\begin{frame}
  \titlepage
\end{frame}

\begin{frame}{Session plan (90 minutes)}
  \begin{enumerate}
    \item Why ERP/CRM/SCM integration is central to commerce
    \item The enterprise systems perspective
    \item Commerce versus enterprise ownership
    \item Systems of record for key entities
    \item Master data management (MDM)
    \item Managerial implications
  \end{enumerate}
\end{frame}

\begin{frame}{Learning objectives}
  By the end of this lecture, you should be able to:
  \begin{itemize}
    \item Explain how commerce platforms integrate with ERP, CRM, and supply-chain systems across operational workflows.
    \item Identify systems of record for key entities such as customer, product, price list, order, invoice, and shipment.
    \item Analyze master data management, process ownership, and consistency risks in hybrid commerce architectures.
    \item Design integration flows for order-to-cash, returns, and B2B pricing using Odoo as a reference ERP environment.
  \end{itemize}
\end{frame}

\section{Why ERP/CRM/SCM integration is central to commerce}

\begin{frame}{Why ERP/CRM/SCM integration is central to commerce}
  \begin{itemize}
    \item orders that appear confirmed online but fail in back-office processing,
    \item inventory figures that differ between storefront and warehouse,
    \item duplicate customer records across channels,
    \item B2B price lists that drift between sales and commerce systems,
    \item returns and refunds that create reconciliation problems in finance.
  \end{itemize}
\end{frame}

\section{The enterprise systems perspective}

\begin{frame}{The enterprise systems perspective}
  \begin{itemize}
    \item \textbf{ERP:} \textbf{Sales:} quotations, sales orders, customer terms, sales workflows,; \textbf{Inventory:} stock moves, reservations, warehouse operations, replenishment,
    \item \textbf{CRM:} CRM systems focus on relationship history, account management, sales interactions, lead qualification, and customer-service context.
    \item \textbf{SCM and warehouse systems:} Supply-chain and warehouse platforms manage sourcing, supplier coordination, inbound logistics, warehouse execution, and delivery orchestration.
  \end{itemize}
\end{frame}

\section{Commerce versus enterprise ownership}

\begin{frame}{Commerce versus enterprise ownership}
  \begin{itemize}
    \item \textbf{Common ownership patterns:} The \textbf{commerce layer} often owns web sessions, carts, checkout state, personalization context, and channel-specific merchandising behavior.; The \textbf{ERP layer} often owns invoices, accounting entries, stock movements, contractual customer records, and internal fulfillment commitments.
    \item \textbf{Why ownership must be explicit:} If pricing is editable in both storefront and ERP, or if both commerce and CRM create customer records independently, inconsistencies are not accidental.
  \end{itemize}
\end{frame}

\section{Systems of record for key entities}

\begin{frame}{Systems of record for key entities}
  \begin{itemize}
    \item \textbf{Customer:} the storefront may own channel account credentials and preferences,; ERP may own bill-to and ship-to parties for financial processing,
    \item \textbf{Product:} ERP may own procurement and internal stock-keeping information,; a catalog or PIM layer may own rich digital merchandising content,
    \item \textbf{Price list:} Price lists are especially sensitive in B2B environments.
    \item \textbf{Order:} The commerce platform often captures the initial order intent, but the enterprise question is whether ERP becomes the authoritative record for downstream operational and financial lifecycle.
    \item \textbf{Invoice and payment reconciliation:} Invoices, ledger effects, and payment reconciliation usually belong in ERP or finance-owned systems.
    \item \textbf{Shipment:} Shipment truth may belong in warehouse or logistics systems even if ERP or storefront mirrors the status.
  \end{itemize}
\end{frame}

\section{Master data management (MDM)}

\begin{frame}{Master data management (MDM)}
  \begin{itemize}
    \item \textbf{Why MDM matters:} duplicate customer records created by multiple channels,; product records with inconsistent category or attribute definitions,
    \item \textbf{Practical MDM principles:} Define authoritative identifiers and mapping rules.; Separate creation rights from consumption rights.
  \end{itemize}
\end{frame}

\section{Order-to-cash in integrated commerce}

\begin{frame}{Order-to-cash in integrated commerce}
  \begin{itemize}
    \item \textbf{Typical flow:} customer places an order in the storefront,; payment is authorized or payment terms are validated,
    \item \textbf{Architectural decisions inside the flow:} whether the storefront can confirm an order before ERP acceptance,; whether stock is reserved in real time or reconciled later,
    \item \textbf{Failure modes:} the order is captured online but fails to post in ERP,; payment succeeds but the order export fails,
  \end{itemize}
\end{frame}

\section{Inventory, fulfillment, and availability}

\begin{frame}{Inventory, fulfillment, and availability}
  \begin{itemize}
    \item \textbf{Inventory concepts that should not be conflated:} \textbf{On-hand inventory:} physical stock currently available in storage.; \textbf{Reserved inventory:} stock allocated to pending commitments.
    \item \textbf{Synchronization strategies:} near-real-time events for stock adjustments,; short-lived availability caches,
  \end{itemize}
\end{frame}

\section{Returns and refunds}

\begin{frame}{Returns and refunds}
  \begin{itemize}
    \item \textbf{Typical return flow:} customer initiates a return,; policy eligibility is checked,
    \item \textbf{Architectural risks:} the storefront approves a return that finance policy rejects,; inventory is restocked before inspection rules are applied,
  \end{itemize}
\end{frame}

\section{B2B pricing and account structures}

\begin{frame}{B2B pricing and account structures}
  \begin{itemize}
    \item \textbf{Price lists and negotiated terms:} customer account,; contract terms,
    \item \textbf{Account hierarchies:} a legal entity,; several branches,
  \end{itemize}
\end{frame}

\section{Using Odoo as a reference ERP}

\begin{frame}{Using Odoo as a reference ERP}
  \begin{itemize}
    \item \textbf{Illustrative Odoo-aligned mapping:} \textbf{Sales:} quotations, orders, commercial terms, B2B account workflows.; \textbf{Inventory:} stock movements, reservations, warehouse updates.
    \item \textbf{What should remain vendor-neutral:} system-of-record analysis,; event and API contracts,
  \end{itemize}
\end{frame}

\section{Integration risks and mitigation strategies}

\begin{frame}{Integration risks and mitigation strategies}
  \begin{itemize}
    \item \textbf{Duplicate sources of truth:} mitigated by explicit ownership and reconciliation.
    \item \textbf{Latency mismatch:} mitigated by choosing appropriate sync versus async boundaries.
    \item \textbf{Identifier mismatch:} mitigated by governed master identifiers and mapping tables.
    \item \textbf{Status divergence:} mitigated by explicit lifecycle models and state-mapping rules.
    \item \textbf{Operational blind spots:} mitigated by observability, error queues, and business-level monitoring.
  \end{itemize}
\end{frame}

\section{Managerial implications}

\begin{frame}{Managerial implications}
  \begin{itemize}
    \item Enterprise integration decisions affect more than system design.
    \item They shape service reliability, customer trust, financial control, and the ability to scale channels.
    \item If commerce and ERP are tightly coupled in every user interaction, customer experience becomes hostage to back-office latency.
  \end{itemize}
\end{frame}

\end{document}
